Finden Sie den Randoperator $\partial$ für das Netzwerk
\begin{center}
\def\p{
\raisebox{-0.2cm}{\begin{tikzpicture}[>=latex,thick]
\draw[color=gray] (0,0) rectangle (0.9,0.7);
\end{tikzpicture}}
}
\begin{tikzpicture}[>=latex,thick]
\def\l{1.3}
\coordinate (A1) at ({-sqrt(3)*\l/2},{0.5*\l});
\coordinate (A2) at ({sqrt(3)*\l/2},{0.5*\l});
\coordinate (A3) at (0,{-\l});
\draw[->,shorten <= 0.25cm,shorten >= 0.25cm] (A1) -- (A2);
\draw[->,shorten <= 0.25cm,shorten >= 0.25cm] (A2) -- (A3);
\draw[->,shorten <= 0.25cm,shorten >= 0.25cm] (A1) -- (A3);
\node at (A1) {$1$}; \draw (A1) circle[radius=0.25];
\node at (A2) {$2$}; \draw (A2) circle[radius=0.25];
\node at (A3) {$3$}; \draw (A3) circle[radius=0.25];
\node at ($0.5*(A1)+0.5*(A2)$) [above] {$1$};
\node at ($0.5*(A1)+0.5*(A3)$) [left] {$2$};
\node at ($0.5*(A2)+0.5*(A3)$) [right] {$3$};
\node at (5,-0.2) {$\displaystyle
\partial =
\renewcommand{\arraystretch}{1.4}
\begin{pmatrix}
\p & \p & \p \\[2pt]
\p & \p & \p \\[2pt]
\p & \p & \p
\end{pmatrix}
$};
\end{tikzpicture}
\end{center}
\begin{loesung}
\[
\partial
=
\begin{pmatrix*}[r]
-1& -1 &  0 \\
 1&  0 & -1 \\
 0&  1 &  1 
\end{pmatrix*}
\qedhere
\]
\end{loesung}
